% Part: incompleteness
% Chapter: incompleteness-provability
% Section: lob-thm

\documentclass[../../../include/open-logic-section]{subfiles}

\begin{document}

\olfileid{inc}{inp}{lob}

\olsection{लोबचे प्रमेय}

उपपत्ती~$\Th{T}$ साठीचे गोडेल-वाक्य
हा $\lnot \OProv[\Th{T}](y)$ चा
स्थिरबिंदू असतो; म्हणजे असे
!!a{sentence}~$!G$ असते की
\[
\Th{T} \Proves \lnot \OProv[\Th{T}](\gn{!G}) \liff !G.
\]
ते !!{derivable} नसते. कारण
$\Th{T} \Proves !G$ असेल, तर
(a) !!{derivability} अट~(1) वरून
$\Th{T} \Proves \OProv[\Th{T}](\gn{!G})$,
आणि (b) $\Th{T}
\Proves !G$ व $\Th{T} \Proves \lnot \OProv[\Th{T}](\gn{!G})
\liff !G$ यांवरून
$\Th{T} \Proves \lnot \OProv[\Th{T}](\gn{!G})$
मिळते. त्यामुळे $\Th{T}$ विसंगत
ठरेल. आता $\OProv[\Th{T}](y)$
च्या स्थिरबिंदूची स्थिती विचारणे
स्वाभाविक आहे. म्हणजे असे
!!a{sentence}~$!H$ घेऊ की
\[
\Th{T} \Proves \OProv[\Th{T}](\gn{!H}) \liff !H.
\]
ते !!{derivable} असेल, तर
अट~(1) वरून
$\Th{T} \Proves \OProv[\Th{T}](\gn{!H})$.
पण वरील सममूल्यतेला
निष्कर्षणाचा नियम लावूनही
हाच निष्कर्ष मिळतो.
म्हणून गोडेल-वाक्याच्या
बाबतीतला युक्तिवाद येथे
$\Th{T}$ विसंगत आहे
असे दाखवत नाही.
अर्थात यावरून
$\Th{T}$ हे
$!H$ ला \emph{खरोखर}
!!{derive} करते असेही
सिद्ध होत नाही.

हा प्रश्न थोडा व्यापक
केल्यास पुढे जाता येते.
स्थिरबिंदूच्या
सममूल्यतेची
डावीकडून उजवीकडची
दिशा
$\OProv[\Th{T}](\gn{!H}) \lif !H$
ही \emph{प्रतिबिंबन तत्त्व}
म्हणवणाऱ्या सर्वसाधारण
रूपबंधाचे एक उदाहरण आहे:
$\OProv[\Th{T}](\gn{!A}) \lif !A$.
त्याला असे नाव दिले
आहे कारण ते एका
अर्थाने $\Th{T}$
स्वतः काय !!{derive}
करू शकते यावर
“प्रतिबिंबन” करू शकते
असे व्यक्त करते.
मूलतः ते कोणत्याही~$!A$
साठी “$\Th{T}$
$!A$ ला !!{derive}
करू शकत असेल, तर
$!A$ सत्य आहे”
असे म्हणते. हे
अर्थात फक्त सत्यनिष्ठ
उपपत्त्यांसाठीच खरे
असते. त्यामुळे
उपपत्त्या सर्वसाधारणपणे
त्याचे प्रत्येक उदाहरण
!!{derive} करणार नाहीत
असे वाटते. मग
पुरेशी प्रबळ आणि
!!{derivability} अटी
पूर्ण करणारी उपपत्ती
कोणती उदाहरणे
!!{derive} करू शकते?
$!A$ स्वतः
!!{derivable} असेल
अशी सर्व उदाहरणे
नक्कीच. पुढील
निकाल दाखवतो की
एवढीच उदाहरणे शक्य
आहेत.

\begin{thm}\ollabel{thm:lob}
$\Th{T}$ ही~$\Th{Q}$
चा विस्तार असलेली
!!a{axiomatizable}
उपपत्ती असू द्या,
आणि $\OProv[\Th{T}](y)$
हे सूत्र \olref[2in]{sec}
मधील P1--P3 अटी
पूर्ण करत असेल.
$\Th{T}$ हे
$\OProv[\Th{T}](\gn{!A}) \lif !A$
!!{derive}s करत
असेल, तर प्रत्यक्षात
$\Th{T}$ हे
$!A$ सुद्धा
!!{derive}s करते.
\end{thm}

दुसऱ्या शब्दांत,
$\Th{T} \Proves/ !A$
असेल, तर
$\Th{T} \Proves/
\OProv[\Th{T}](\gn{!A}) \lif !A$.
हा निकाल लोबचे
प्रमेय म्हणून ओळखला
जातो.

\begin{explain}
लोबच्या प्रमेयामागील
युक्ती समजण्यासाठी
सांताक्लॉज अस्तित्वात
आहे हे “सिद्ध”
करणारा चतुर युक्तिवाद
पाहा. (हा निष्कर्ष
आवडत नसेल, तर
तुम्हाला हवा असलेला
दुसरा निष्कर्ष
त्याच्या जागी
घेता येईल.)
युक्तिवाद असा:
\begin{enumerate}
\item $X$ हे “$X$
  सत्य असेल, तर
  सांताक्लॉज
  अस्तित्वात आहे”
  हे वाक्य मानू.
\item $X$ सत्य आहे
  असे धरू.
\item मग त्यात
  म्हटलेले खरे आहे:
  $X$ सत्य असेल,
  तर सांताक्लॉज
  अस्तित्वात आहे.
\item आपण $X$ सत्य
  आहे असे धरलेच
  आहे; म्हणून
  (2) आणि~(3)
  वर निष्कर्षणाचा
  नियम लावून
  सांताक्लॉज
  अस्तित्वात आहे
  असा निष्कर्ष
  काढू शकतो.
\item (2) मधील
  “$X$ सत्य आहे”
  या गृहीतकावरून
  (4) मधील
  “सांताक्लॉज
  अस्तित्वात आहे”
  हे आपण निष्पन्न
  केले. अटीची
  सिद्धता वापरून
  “$X$ सत्य असेल,
  तर सांताक्लॉज
  अस्तित्वात आहे”
  हे दाखवले.
\item पण हेच
  वाक्य~$X$ आहे.
  म्हणजे आपण
  $X$ सत्य आहे
  हे दाखवले.
\item मग वरील
  (2)--(4) च्या
  युक्तिवादाने
  सांताक्लॉज
  अस्तित्वात आहे.
\end{enumerate}
या कल्पनेचे
आकारिकीकरण करताना
“सत्य आहे” ऐवजी
“!!{derivable} आहे”
आणि “सांताक्लॉज
अस्तित्वात आहे”
ऐवजी~$!A$
घेतल्यास लोबच्या
प्रमेयाची सिद्धता
मिळते. युक्ती अशी:
स्थिरबिंदू उपपत्ती
!!{formula}~
$\OProv[\Th{T}](y) \lif !A$
याला लावायची.
त्याचा स्थिरबिंदू
वरील रूपरेषेतील
वाक्य~$X$
याच्याशी जुळतो.
\end{explain}

\begin{proof}[याची सिद्धता: \olref{thm:lob}]
$!A$ हे असे
!!a{sentence}
असू द्या की
$\Th{T}$ हे
$\OProv[\Th{T}](\gn{!A}) \lif !A$
!!{derive}s करते.
$!B(y)$ हे
!!{formula}~
$\OProv[\Th{T}](y)
\lif !A$ मानू.
स्थिरबिंदू उपपत्ती
वापरून असे
!!a{sentence}~$!D$
शोधू की
$\Th{T}$ हे
$!D \liff !B(\gn{!D})$
!!{derive}s करते.
मग पुढीलपैकी
प्रत्येक विधान
$\Th{T}$ मध्ये
!!{derivable} आहे:
\begin{align}
  & !D \liff (\OProv[\Th{T}](\gn{!D}) \lif !A) \ollabel{L-1}\\
  & \qquad \text{$!D$ हा~$!B(y)$ चा स्थिरबिंदू आहे}\notag \\
  & !D \lif (\OProv[\Th{T}](\gn{!D}) \lif !A) \ollabel{L-2}\\
  & \qquad\text{\olref{L-1} वरून}\notag\\
  & \OProv[\Th{T}](\gn{!D \lif (\OProv[\Th{T}](\gn{!D}) \lif !A)}) \ollabel{L-3}\\
  & \qquad \text{\olref{L-2} आणि अट P1 वरून}\notag \\
  & \OProv[\Th{T}](\gn{!D}) \lif \OProv[\Th{T}](\gn{\OProv[\Th{T}](\gn{!D}) \lif !A})
  \ollabel{L-4}\\
  &\qquad \text{\olref{L-3} व अट P2 वरून}\notag \\
  & \OProv[\Th{T}](\gn{!D}) \lif (\OProv[\Th{T}](\gn{\OProv[\Th{T}](\gn{!D})}) \lif \OProv[\Th{T}](\gn{!A})) \ollabel{L-5}\\
  &\qquad \text{\olref{L-4} व अट P2 पुन्हा वापरून} \notag\\
& \OProv[\Th{T}](\gn{!D}) \lif \OProv[\Th{T}](\gn{\OProv[\Th{T}](\gn{!D})}) \ollabel{L-6}\\
  & \qquad\text{!!{derivability} अट P3 वरून} \notag\\
  & \OProv[\Th{T}](\gn{!D}) \lif \OProv[\Th{T}](\gn{!A}) \ollabel{L-7} \\
  &\qquad\text{\olref{L-5} आणि \olref{L-6} वरून}\notag\\
  & \OProv[\Th{T}](\gn{!A}) \lif !A \ollabel{L-8}\\
  &\qquad\text{प्रमेयाच्या गृहीतकावरून} \notag\\
  & \OProv[\Th{T}](\gn{!D}) \lif !A \ollabel{L-9}\\
  &\qquad\text{\olref{L-7} आणि \olref{L-8} वरून}\notag\\
  & (\OProv[\Th{T}](\gn{!D}) \lif !A) \lif !D \ollabel{L-10}\\
  & \qquad \text{\olref{L-1} वरून}\notag \\
  & !D \ollabel{L-11}\\
  & \qquad\text{\olref{L-9} आणि \olref{L-10} वरून}\notag \\
  & \OProv[\Th{T}](\gn{!D}) \ollabel{L-12}\\
  & \qquad\text{\olref{L-11} आणि अट~P1 वरून}\notag \\
  & !A \qquad\qquad\text{\olref{L-8} आणि \olref{L-12} वरून}\notag
\end{align}
\end{proof}

लोबचे प्रमेय मिळाल्यावर
दुसऱ्या अपूर्णता प्रमेयाची
लहान सिद्धता मिळते
(P1--P3 अटी पूर्ण
करणारे !!a{derivability}
विधेय असलेल्या उपपत्त्यांसाठी):
$\Th{T} \Proves
\OProv[\Th{T}](\gn{\lfalse}) \lif \lfalse$
असेल, तर
$\Th{T} \Proves \lfalse$.
$\Th{T}$ सुसंगत
असेल, तर
$\Th{T} \Proves/ \lfalse$.
म्हणून
$\Th{T}
\Proves/ \OProv[\Th{T}](\gn{\lfalse}) \lif \lfalse$,
म्हणजे
$\Th{T} \Proves/
\OCon[\Th{T}]$.
$\OProv[\Th{T}](x)$
चा स्थिरबिंदू~$!H$
हा !!{derivable}
आहे हेही यावरून
दाखवता येते.
कारण
\begin{align*}
  \Th{T} & \Proves \OProv[\Th{T}](\gn{!H}) \liff !H\\
  \intertext{यावरून विशेषतः}
    \Th{T} & \Proves \OProv[\Th{T}](\gn{!H}) \lif !H
\end{align*}
आणि म्हणून लोबच्या
प्रमेयावरून
$\Th{T} \Proves !H$.

% उलट दिशेने, गृहपाठात दुसरे अपूर्णता प्रमेय वापरून
% स्थिरबिंदू उपपत्तीशिवाय लोबच्या प्रमेयाची छोटी सिद्धता
% काढण्यास सांगता येईल.

\begin{prob}
$\Th{T}$ ही संगणनक्षमपणे
स्वयंसिद्धकीकृत उपपत्ती
असू द्या, आणि
$\OProv[\Th{T}]$ हे
$\Th{T}$ साठीचे
!!a{derivability}
विधेय असू द्या.
पुढील चार विधाने
विचारात घ्या:
\begin{enumerate}
\item जर $T \Proves !A$, तर $T \Proves \OProv[\Th{T}](\gn{!A})$.
\item $T \Proves !A \lif \OProv[\Th{T}](\gn{!A})$.
\item जर $T \Proves \OProv[\Th{T}](\gn{!A})$, तर $T \Proves !A$.
\item $T \Proves \OProv[\Th{T}](\gn{!A}) \lif !A$
\end{enumerate}
ही प्रत्येक विधाने
कोणत्या अटींवर
खरी असतात?
\end{prob}

\end{document}
